\documentclass[10pt,a4paper]{amsart}
\usepackage{amsmath,amssymb,amsthm,graphicx}
\usepackage{verbatim}
\newenvironment{reference}{\comment}{\endcomment}
%\newenvironment{reference}{}{}
\newenvironment{slogan}{\comment}{\endcomment}
\newenvironment{history}{\comment}{\endcomment}

% For commutative diagrams we use Xy-pic
\usepackage[all]{xy}

% We use 2cell for 2-commutative diagrams.
\xyoption{2cell}
% Theorem environments.
%
\theoremstyle{plain}
\newtheorem{theorem}[subsection]{Theorem}
\newtheorem{proposition}[subsection]{Proposition}
\newtheorem{lemma}[subsection]{Lemma}

\theoremstyle{definition}
\newtheorem{definition}[subsection]{Definition}
\newtheorem{example}[subsection]{Example}
\newtheorem{exercise}[subsection]{Exercise}
\newtheorem{situation}[subsection]{Situation}

\theoremstyle{remark}
\newtheorem{remark}[subsection]{Remark}
\newtheorem{remarks}[subsection]{Remarks}

\numberwithin{equation}{subsection}

% Macros
%
\def\lim{\mathop{\mathrm{lim}}\nolimits}
\def\colim{\mathop{\mathrm{colim}}\nolimits}
\def\Spec{\mathop{\mathrm{Spec}}}
\def\Hom{\mathop{\mathrm{Hom}}\nolimits}
\def\Ext{\mathop{\mathrm{Ext}}\nolimits}
\def\SheafHom{\mathop{\mathcal{H}\!\mathit{om}}\nolimits}
\def\SheafExt{\mathop{\mathcal{E}\!\mathit{xt}}\nolimits}
\def\Sch{\mathit{Sch}}
\def\Mor{\mathop{\mathrm{Mor}}\nolimits}
\def\Ob{\mathop{\mathrm{Ob}}\nolimits}
\def\Sh{\mathop{\mathit{Sh}}\nolimits}
\def\NL{\mathop{N\!L}\nolimits}
\def\CH{\mathop{\mathrm{CH}}\nolimits}
\def\proetale{{pro\text{-}\acute{e}tale}}
\def\etale{{\acute{e}tale}}
\def\QCoh{\mathit{QCoh}}
\def\Ker{\mathop{\mathrm{Ker}}}
\def\Im{\mathop{\mathrm{Im}}}
\def\Coker{\mathop{\mathrm{Coker}}}
\def\Coim{\mathop{\mathrm{Coim}}}

% Boxtimes
%
\DeclareMathSymbol{\boxtimes}{\mathbin}{AMSa}{"02}

%
% Macros for moduli stacks/spaces
%
\def\QCohstack{\mathcal{QC}\!\mathit{oh}}
\def\Cohstack{\mathcal{C}\!\mathit{oh}}
\def\Spacesstack{\mathcal{S}\!\mathit{paces}}
\def\Quotfunctor{\mathrm{Quot}}
\def\Hilbfunctor{\mathrm{Hilb}}
\def\Curvesstack{\mathcal{C}\!\mathit{urves}}
\def\Polarizedstack{\mathcal{P}\!\mathit{olarized}}
\def\Complexesstack{\mathcal{C}\!\mathit{omplexes}}
% \Pic is the operator that assigns to X its picard group, usage \Pic(X)
% \Picardstack_{X/B} denotes the Picard stack of X over B
% \Picardfunctor_{X/B} denotes the Picard functor of X over B
\def\Pic{\mathop{\mathrm{Pic}}\nolimits}
\def\Picardstack{\mathcal{P}\!\mathit{ic}}
\def\Picardfunctor{\mathrm{Pic}}
\def\Deformationcategory{\mathcal{D}\!\mathit{ef}}
\usepackage{cleveref}
\crefname{originallemma}{original lemma}{original lemmas}
\Crefname{originallemma}{Original lemma}{Original lemmas}
\makeatletter
\newlabel{algebra-lemma-ring-with-only-minimal-primes}{{04MG}{source}{}{sourcepointer-1}{}}
\newlabel{algebra-lemma-ring-with-only-minimal-primes@cref}{{[originallemma][0][]04MG}{[1][0][]source}{}{}{}}
\makeatother
\crefname{originallemma}{original lemma}{original lemmas}
\Crefname{originallemma}{Original lemma}{Original lemmas}
\makeatletter
\newlabel{schemes-lemma-characterize-quasi-separated}{{01KO}{source}{}{sourcepointer-2}{}}
\newlabel{schemes-lemma-characterize-quasi-separated@cref}{{[originallemma][0][]01KO}{[1][0][]source}{}{}{}}
\makeatother
\crefname{originallemma}{original lemma}{original lemmas}
\Crefname{originallemma}{Original lemma}{Original lemmas}
\makeatletter
\newlabel{schemes-lemma-characterize-separated}{{01KP}{source}{}{sourcepointer-3}{}}
\newlabel{schemes-lemma-characterize-separated@cref}{{[originallemma][0][]01KP}{[1][0][]source}{}{}{}}
\makeatother
\crefname{originallemma}{original lemma}{original lemmas}
\Crefname{originallemma}{Original lemma}{Original lemmas}
\makeatletter
\newlabel{topology-lemma-quasi-compact-in-Hausdorff}{{08YB}{source}{}{sourcepointer-4}{}}
\newlabel{topology-lemma-quasi-compact-in-Hausdorff@cref}{{[originallemma][0][]08YB}{[1][0][]source}{}{}{}}
\makeatother
\title{Characterizations of schemes of dimension zero}
\author{The Stacks Project Authors}
\date{stacks-project:properties.tex:local-snapshot-20261012; GNU FDL 1.2 or later}
\begin{document}
\maketitle
\noindent\textit{Editorial note.} All delivered mathematical text is the author's own, copied from the pinned Stacks Project chapter file properties.tex; no statement or proof step was rewritten, reordered or removed.\par\medskip
\noindent\textit{Editorial note.} The counted proof is the author's complete argument: the cycle of implications through the criteria on affine open coverings, and the two closing deductions: connected components are points via separation by an affine open, and quasi-compact opens are affine by induction using the three-piece decomposition.\par\medskip
\noindent\textit{Editorial note.} The label lemma-dimension-zero names the counted statement, and the proof environment immediately following it in the source is delivered verbatim as its proof.\par\medskip
\noindent\textit{Editorial note.} The counted statement is carried by the author's own lemma environment in the chapter source.\par\medskip
\noindent\textit{Editorial note.} Pre-open gate: python operations/agent-authoring/tools/precard\_gate.py --source properties.tex --label lemma-dimension-zero --start 1519 --end 1587 --local-id 28031 returned RESULT=OK (L28031 free, occupancy FREE, public interval clear). Candidate taken from sources/stacks-properties/candidates-halfA.json (third exclusive pool assigned by the Lead on 2026-10-13).\par\medskip
\noindent\textit{Editorial note.} The statement carries a native bibliography pointer environment naming the source of the argument; that environment is part of the delivered statement and is covered by the preamble span 8-13 of the chapter, which this card includes.\par\medskip
\noindent\textit{Editorial note.} The statement and proof are copied from the pinned chapter file; the cited inputs are kept as external pointers with their canonical Stacks tags.\par\medskip
\noindent\textit{Allowed standard prerequisites (editorial).}\begin{itemize}
\item Structure of the spectrum of a ring with only minimal primes. Algebra, the lemma describing the spectrum of a ring whose primes are all minimal as a profinite space with the properties used freely in the proof. Supplies the properties of an affine open of a scheme of dimension zero, in particular its profiniteness and Hausdorffness.
\item Criteria for quasi-separatedness and separatedness in terms of intersections of affine opens. Schemes, the lemmas characterizing quasi-separated and separated schemes through the intersections of the members of an affine open covering. Supplies the two characterizations applied in the first two implications.
\item Quasi-compact subsets of a Hausdorff space are closed. Topology, the lemma that a quasi-compact subset of a Hausdorff space is closed. Supplies the implication from Hausdorffness to the closedness of affine opens.
\end{itemize}\medskip
\ifdefined\hypertarget\hypertarget{sourcepointer-1}{}\fi
\noindent\textit{Source reference.} Original lemma 04MG: sources/stacks-chow/tags.txt, source lines 6042--6042. Prerequisite pointer invoked at proof line(s) 1540; it names a standard fact of a sibling chapter that the counted argument uses, and every step of the counted proof built on it is retained verbatim in the delivered body.\par\medskip
\ifdefined\hypertarget\hypertarget{sourcepointer-2}{}\fi
\noindent\textit{Source reference.} Original lemma 01KO: sources/stacks-chow/tags.txt, source lines 2057--2057. Prerequisite pointer invoked at proof line(s) 1546; it names a standard fact of a sibling chapter that the counted argument uses, and every step of the counted proof built on it is retained verbatim in the delivered body.\par\medskip
\ifdefined\hypertarget\hypertarget{sourcepointer-3}{}\fi
\noindent\textit{Source reference.} Original lemma 01KP: sources/stacks-chow/tags.txt, source lines 2058--2058. Prerequisite pointer invoked at proof line(s) 1556; it names a standard fact of a sibling chapter that the counted argument uses, and every step of the counted proof built on it is retained verbatim in the delivered body.\par\medskip
\ifdefined\hypertarget\hypertarget{sourcepointer-4}{}\fi
\noindent\textit{Source reference.} Original lemma 08YB: sources/stacks-chow/tags.txt, source lines 11334--11334. Prerequisite pointer invoked at proof line(s) 1571; it names a standard fact of a sibling chapter that the counted argument uses, and every step of the counted proof built on it is retained verbatim in the delivered body.\par\medskip
\begin{lemma}
\label{lemma-dimension-zero}
\begin{reference}
Email from Ofer Gabber dated June 4, 2016
\end{reference}
Let $X$ be a scheme of dimension zero. The following are equivalent
\begin{enumerate}
\item $X$ is quasi-separated,
\item $X$ is separated,
\item $X$ is Hausdorff,
\item every affine open is closed.
\end{enumerate}
In this case the connected components of $X$ are points and every
quasi-compact open of $X$ is affine. In particular, if $X$
is quasi-compact, then $X$ is affine.
\end{lemma}

\begin{proof}
As the dimension of $X$ is zero, we see that for any affine open
$U \subset X$ the space $U$ is profinite and satisfies
a bunch of other properties which we will use freely below, see
Algebra, Lemma \ref{algebra-lemma-ring-with-only-minimal-primes}.
We choose an affine open covering $X = \bigcup U_i$.

\medskip\noindent
If (4) holds, then $U_i \cap U_j$ is a closed subset of
$U_i$, hence quasi-compact, hence $X$ is quasi-separated,
by Schemes, Lemma \ref{schemes-lemma-characterize-quasi-separated},
hence (1) holds.

\medskip\noindent
If (1) holds, then $U_i \cap U_j$ is a quasi-compact open
of $U_i$ hence closed in $U_i$. Then $U_i \cap U_j \to U_i$
is an open immersion whose image is closed, hence it is a
closed immersion. In particular $U_i \cap U_j$ is affine
and $\mathcal{O}(U_i) \to \mathcal{O}_X(U_i \cap U_j)$ is surjective.
Thus $X$ is separated
by Schemes, Lemma \ref{schemes-lemma-characterize-separated},
hence (2) holds.

\medskip\noindent
Assume (2) and let $x, y \in X$. Say $x \in U_i$. If $y \in U_i$
too, then we can find disjoint open neighbourhoods of $x$ and $y$
because $U_i$ is Hausdorff. Say $y \not \in U_i$ and $y \in U_j$.
Then $y \not \in U_i \cap U_j$ which is an affine open of $U_j$
and hence closed in $U_j$. Thus we can find an open neighbourhood
of $y$ not meeting $U_i$ and we conclude that $X$ is Hausdorff,
hence (3) holds.

\medskip\noindent
Assume (3). Let $U \subset X$ be affine open.
Then $U$ is closed in $X$ by Topology, Lemma
\ref{topology-lemma-quasi-compact-in-Hausdorff}.
This proves (4) holds.

\medskip\noindent
Assume $X$ satisfies the equivalent conditions (1) -- (4).
We prove the final statements of the lemma. Say $x, y \in X$
with $x \not = y$. Since $y$ does not specialize to $x$ we can
choose $U \subset X$ affine open with $x \in U$ and $y \not \in U$.
Then we see that $X = U \amalg (X \setminus U)$ is a decomposistion
into open and closed subsets which shows that $x$ and $y$ do not belong
to the same connected component of $X$. Next, assume $U \subset X$
is a quasi-compact open. Write $U = U_1 \cup \ldots \cup U_n$ as a union
of affine opens. We will prove by induction on $n$ that $U$ is affine.
This immediately reduces us to the case $n = 2$. In this case we have
$U = (U_1 \setminus U_2) \amalg (U_1 \cap U_2) \amalg (U_2 \setminus U_1)$
and the arguments above show that each of the pieces is affine.
\end{proof}

\end{document}
